\section{Sharp inference in every degree}
\label{sec:polynomial-statistics}
The observable inverse gives inference without a latent channel floor.
We first prove sharp signal-class rates with simple component roots,
and then prove sharp multiplicity-dependent rates on neighbourhoods of
polynomial singularities. The latter are neighbourhood minimax rates,
not a pointwise lower bound at a single fixed parameter.

Take $n$ independent categorical draws at each of the $\ell$ fixed
clocks, put $N=\ell n$, $r=m_1m_2$, and define
\begin{equation}\label{eq:poly-confidence}
 \rho=\sqrt{\frac{r}{2n}\log\frac{2\ell r}{\zeta}},\qquad
 \mathcal D=\{\omega\in\cK_d:
       \max_j\|P_\omega(T_j)-\widehat P_j\|_F\le\rho\}.
\end{equation}
Use the same empty-set convention as in \eqref{eq:confidence}. Write
$\mathcal D_\Lambda$ for its multiset image.

\begin{theorem}[Observable polynomial inference]\label{thm:poly-inference}
The set \eqref{eq:poly-confidence} has coverage at least $1-\zeta$
uniformly over the entire closed model. If the true datum has
$\eta_d(P)>0$, the direct estimator satisfies
\begin{align}
 \E_Pd_\infty(\widehat\Lambda,\Lambda)^2
 &\le C\min\{1,(N\eta_d(P)^2)^{-1/d}\},\nonumber\\
 \E_P\diam_\infty\mathcal D_\Lambda
 &\le C_\zeta\min\{1,(\sqrt N\eta_d(P))^{-1/d}\}.
                                                       \label{eq:poly-upper}
\end{align}
The exponent $1/d$ is replaced by $1/m$ under the separated-cluster
hypothesis in Theorem~\ref{thm:joint}. In particular, simple component
roots separated by $\nu$ give squared risk
$C_\nu\min\{1,(N\eta_d(P)^2)^{-1}\}$ and expected diameter
$C_{\nu,\zeta}\min\{1,(\sqrt N\eta_d(P))^{-1}\}$.
\end{theorem}
\begin{proof}
The union bound for the $\ell r$ cell frequencies gives
$\Pr(e>\rho)\le2\ell r\exp(-2n\rho^2/r)=\zeta$.
Also $\E e^2\le\sum_j\E\|\widehat P_j-P_j\|_F^2\le\ell/n=\ell^2/N$.
Apply Theorem~\ref{thm:joint} to the direct estimator, and use Jensen's
inequality for the power $1/d$ of $e^2$. For nonempty residual regions,
Lemma~\ref{lem:modulus} bounds their diameter by
$C\min\{1,[(e+\rho)/\eta_d(P)]^{1/d}\}$ on every outcome. Taking
expectations gives the second assertion; the empty-set convention
preserves it. The identical argument proves the cluster refinements.
\end{proof}

\begin{theorem}[Sharp semisimple signal classes]\label{thm:poly-sharp}
Fix $d,k,m_1,m_2$ and $\ell\ge2d+1$. There exist $\nu_0,s_0>0$ such
that, for every fixed $0<\nu\le\nu_0$ and $0<s\le s_0$, the class
\[
 \mathcal S_{\nu,s}=\{\omega\in\cK_d:
        \eta_d(P_\omega)\ge s,\quad
        \hbox{roots within each }f_b\hbox{ are }\nu\hbox{-separated}\}
\]
satisfies
\begin{equation}\label{eq:poly-minimax}
 \inf_{\widehat\Lambda}\sup_{\omega\in\mathcal S_{\nu,s}}
       \E_\omega d_\infty(\widehat\Lambda,\Lambda)^2
           \asymp\min\{1,(Ns^2)^{-1}\}.
\end{equation}
The balanced fixed design attains the upper bound. The lower bound
holds even if the infimum includes all parameter-independent adaptive
clock strategies on $[T_1,T_\ell]$. For confidence sets honest at level
$1-\zeta$, $0<\zeta<1/2$, over the closed model, the corresponding
worst-case expected diameter on this class has order
$\min\{1,(\sqrt Ns)^{-1}\}$.
\end{theorem}
\begin{proof}
Only the lower bound remains. Choose a monic $F$ with $d$ simple
interior roots and a sufficiently small nonzero constant $G$. Choose
$t_0>0$ so that, for $0\le t\le t_0$, $\lambda=1-t$ and the polynomials
\begin{equation}\label{eq:degree-product}
 f_{1,t}=F+\lambda G,\qquad f_{2,t}=F-\lambda G,
 \qquad f_{b,t}=F\quad(b\ge3)
\end{equation}
have simple interior roots with a common positive within-component
separation. Their collection is relatively prime and, with equal
weights, $q=F$. Use the positive stochastic channels of
\eqref{eq:product-family}, now with this $\lambda$. The extra directions
are unchanged, and take sufficiently small contrasts
$\epsilon_u,\epsilon_v>0$. Put $w=\epsilon_u\epsilon_v$.
Direct expansion gives
\begin{align}
 P_t(T)&=A_t+\frac{G}{F(T)}D,\qquad
 D=\frac{\epsilon_u ev_0^{\mathsf T}+\epsilon_vu_0f^{\mathsf T}}{k},
                                                        \label{eq:degree-identity}\\
 P_t(T)-P_0(T)&=\frac{w}{2k}(\lambda^{-2}-1)ef^{\mathsf T}.
                                                        \label{eq:degree-difference}
\end{align}
Both observed marginals are fixed, and every cell has a uniform positive
floor on the whole clock interval. The leading matrix is
$A_t=U_tV_t^{\mathsf T}/k$, so the orthogonal pair-direction calculation
in Lemma~\ref{lem:families} gives
$\sigma_k(A_t)=w/(2k\lambda^2)$ and $\kappa=w/(2\lambda^2)$.
Consequently \eqref{eq:joint-comparison} gives $B_d(P_t)\asymp w^{-1}$.

The normalization operator in this family is the tensor product of
$r\mapsto\Omega(r(T_j)G/F(T_j))_j$ and $\operatorname{vec}D$.
The scalar map is injective: an interpolating polynomial would give
$FH-rG=0$ at $2d+1$ nodes, hence identically, and coprimality forces
$F\mid r$, so $r=0$. Its least singular value is a fixed positive
number. The two summands in $D$ are Frobenius orthogonal; hence
$\tau(P_t)\asymp\max\{\epsilon_u,\epsilon_v\}$. Uniformly over
the declared ranges, therefore,
\begin{equation}\label{eq:degree-signal}
 c_1w\le\eta_d(P_t)\le C_1w,\qquad
 \sup_T\|P_t(T)-P_0(T)\|_F\le C_2wt.
\end{equation}
No common polynomial factor was attached to an affine experiment:
$\gcd(F,G)=1$ and the visible denominator has degree $d$.

The roots of $F\pm\lambda G$ have nonzero derivative with respect to
$\lambda$, since their derivatives are $\mp G/(F')$ at these roots.
Shrinking $t_0$ preserves their order relative to the roots of $F$ and
gives $c_3t\le d_\infty(\Lambda_t,\Lambda_0)\le C_3t$.
All constants here are fixed before choosing $N,w$.
The positive cell floor and \eqref{eq:degree-signal} imply
$\sup_T\KL(P_t(T),P_0(T))\le Cw^2t^2$. The adaptive transcript
chain rule from Theorem~\ref{thm:statistics} gives the bound
$CNw^2t^2$ for every allowed policy. Take
$t=c\min\{t_0,(\sqrt Nw)^{-1}\}$ to keep total variation bounded
away from one. The two-point risk is at least $ct^2$.

For the exact signal class choose $M\ge1/c_1$, set $w=Ms$ and
$\epsilon_u=\epsilon_v=\sqrt{Ms}$, and restrict $s_0$ so that all
parameters remain in their fixed admissible ranges. Equation
\eqref{eq:degree-signal} puts both testing points, and the path between
them, in $\mathcal S_{\nu,s}$. Choose $\nu_0$ from the separation of
\eqref{eq:degree-product}. This proves \eqref{eq:poly-minimax}.
For honest sets, transfer coverage from the alternative to the base
law and use
$\Pr_0\{\Lambda_0,\Lambda_t\in J\}\ge1-2\zeta-\TV$.
Choosing $c$ in terms of $\zeta$ proves the diameter lower bound.
Theorem~\ref{thm:poly-inference} proves the matching upper bound.
\end{proof}

\subsection{Neighbourhoods of multiple component roots}
The exponent $1/d$ is not merely an artefact of allowing complex roots.
Real-rooted alternatives can force its multiplicity-refined version.
The correct lower-bound construction uses two moving alternatives,
not an arbitrary constant perturbation of $(z-x)^m$.

\begin{theorem}[Multiplicity-dependent local minimax exponents]
\label{thm:multiplicity}
Let $\omega_0$ have strictly positive full-rank channels, strict weights,
$\tau(P_0)>0$, and pairwise distinct component polynomials whose roots
are in $(L,D)$. Let $m$ be the largest multiplicity of a root within a
single component. There is a sufficiently small fixed neighbourhood
$\mathcal U$ of $\omega_0$, in channel, weight and monic-coefficient
coordinates and restricted to the real-rooted model, such that
\begin{equation}\label{eq:local-minimax}
 \inf_{\widehat\Lambda}\sup_{\omega\in\mathcal U}
     \E_\omega d_\infty(\widehat\Lambda,\Lambda)^2
       \asymp N^{-1/m}\qquad(N\longrightarrow\infty).
\end{equation}
The direct estimator at the balanced design attains the upper bound.
The lower bound remains valid for all adaptive clock strategies on
$[T_1,T_\ell]$. The corresponding worst-case expected diameter of
sets honest on $\mathcal U$ has lower order $N^{-1/(2m)}$, attained
by \eqref{eq:poly-confidence}, which is honest even on the closed model.
\end{theorem}
\begin{proof}
Choose disjoint small intervals about the distinct roots of each base
component polynomial. By continuity of roots, a sufficiently small
coefficient neighbourhood has the same number of roots in each such
interval. There are at most $m$ roots per component interval; distances
between its different intervals have a positive lower bound. Keep the
$k$ complete polynomial tuples distinct and the channels full rank by
shrinking the neighbourhood. The joint projectors, $B_d$ and $\tau$
are continuous there; in particular $\eta_d\ge s_*>0$. The cluster
version of Theorem~\ref{thm:joint} is uniform on this neighbourhood,
even when roots inside a cluster split or collide. Theorem
\ref{thm:poly-inference} gives the required risk and diameter upper
bounds with exponent $1/m$.

Choose one base component $f_b(z)=(z-x)^m a(z)$ with $a(x)\ne0$.
Let $H(y)$ be a monic degree-$m$ polynomial with distinct real roots.
A sufficiently small nonzero constant $c$ makes both $H$ and $H+c$
have $m$ distinct real roots, by disjoint sign-changing brackets.
For small $u>0$ replace only this component by the two alternatives
\begin{equation}\label{eq:moving-alternatives}
 f_{b,u,0}(z)=u^mH((z-x)/u)a(z),\qquad
 f_{b,u,1}(z)=u^m\{H((z-x)/u)+c\}a(z).
\end{equation}
All other parameters stay fixed. These are monic, real-rooted and
converge in coefficients to the original component. Thus both lie in
$\mathcal U$ for sufficiently small $u$, and their difference is
exactly $cu^m a(z)$. Their aggregate root multisets are at distance
$c_0u$ for some $c_0>0$. Indeed the changed cluster is the dilation by
$u$ of two different fixed multisets; unchanged roots at $x$, if any,
add the same finite number of zeros to both and cannot make those
multisets equal. Distinct base clusters remain separate.

On the compact clock interval the denominator has a positive lower
bound, so the rational observations differ by at most $Cu^m$ uniformly
in the clock. The positive channels give a common cell floor. Hence
every adaptive transcript has relative entropy at most $CNu^{2m}$.
Choose $u=c'N^{-1/(2m)}$. The two-point test gives squared risk at
least $c''u^2$, proving \eqref{eq:local-minimax}. Coverage transfer
between these same two moving alternatives gives expected diameter
at least a fixed multiple of $u$. This proves the final assertion.
\end{proof}

The local theorem is valid for every $1\le m\le d$ for which such a
regular-normalization base exists; for instance one component may have
an $m$-fold root and another a polynomial with no common zero with it.
Cross-component multiplicity alone does not change $m$. The sharp
neighbourhood rate is therefore compatible with linear semisimple
stability. Intersections with $\tau=0$ or channel-rank loss require
their own signal geometry; the global certificate, the sharp
semisimple classes, the retained affine collapse family, and the local
multiplicity theorem state separately the conclusions proved here.
